
\chapter{Análise do Código MATLAB: \texttt{tri2link.m}}
\section{Visão Geral}

\textbf{What is the main purpose of this file?} \\
The main purpose of this script is to convert a 3D triangular mesh (represented by face indices $T$ and vertex coordinates $V$) into an ``anchored link'' data structure. This structure consists of geometric edge information ($g$), an edge-connecting permutation map ($\phi$), and a source vertex map ($s$).

\textbf{What type of analysis or calculation does it perform?} \\
It performs topological indexing and discrete geometric calculations. Specifically, it maps out the connectivity of directed edges within each triangle and calculates the logarithmic ``amplitwist''---the natural logarithm of the length ratio and the angle difference (twist) between adjacent directed edges along a face.

\textbf{What is the mathematical/theoretical context?} \\
The code operates within the context of \textbf{discrete differential geometry}, \textbf{algebraic topology} (specifically simplicial complexes), and \textbf{graph theory} (directed graphs embedded on surfaces). Calculating ratios and angles to form a complex variable ($g = \log(r) + i\theta$) strongly relates to discrete conformal geometry and discrete complex analysis.

\section{Análise Passo a Passo do Código}

\subsection{1. Setup and Variable Initialization}
\textbf{Explanation:} \\
This section initializes the necessary utility functions and extracts the dimensions of the input matrix. It creates an anonymous function \texttt{id} that generates an identity index map (1 to $N$) for any given vector. It then extracts the dimensions of the triangle matrix $T$, determining the number of faces (\texttt{nT}) and the number of vertices/edges per face (\texttt{nf}, which is inherently 3 for triangles).

\begin{lstlisting}[language=Matlab]
id=@(x) (1:numel(x))';    
[nT,nf]=size(T);   
\end{lstlisting}

\subsection{2. Topological Mapping ($s$ and $\phi$)}
\textbf{Explanation:} \\
This block constructs the topological connectivity of the mesh. First, it transposes and flattens the triangle matrix $T$ into a column vector $s$, which effectively creates a list of source vertices for every directed edge in the mesh. Next, it creates $\phi$, a permutation map for the edges. It reshapes a linear index array, shifts the rows by one position (moving the first edge to the second, second to third, and third to first), and flattens it. The result is a mapping where each edge points to the next consecutive edge within the same triangular face.

\begin{lstlisting}[language=Matlab]
s=T'; 
s=s(:); 
phi=reshape(1:nf*nT,nf,nT);
phi=phi([2:nf 1],:);
phi=phi(:); 
\end{lstlisting}

\subsection{3. Geometric Calculations ($g$)}
\textbf{Explanation:} \\
This section calculates the discrete geometric properties. It determines the target vertex $t$ for each edge using the $\phi$ map. It then defines anonymous functions to compute 3D direction vectors (\texttt{dir}) and their Euclidean norms (\texttt{norm}). 

Using these, it calculates the \texttt{ampli} (the ratio of the length of the next edge to the current edge) and the \texttt{twist} (the angle between the current edge and the next edge, calculated using the dot product and arccosine). Finally, it constructs the complex vector $g$ using the formula $g = \log(\text{ampli}) + 1i \cdot \text{twist}$.

\begin{lstlisting}[language=Matlab]
t=s(phi); idA=id(s);
dir=@(a) V(t(a),:)-V(s(a),:); 
norm=@(v) sqrt(dot(v',v')');   
ampli=norm(dir(phi))./norm(dir(idA));
twist=acos(dot(dir(phi)',dir(idA)')'./(norm(dir(phi)).*norm(dir(idA))));
g=log(ampli)+1i*twist;
\end{lstlisting}

\section{Saídas e Elementos Visuais Esperados}

\textbf{Explanation:} \\
The \texttt{tri2link} function strictly performs mathematical transformations and data restructuring; it \textbf{does not generate any graphical plots, figures, or console visualizations natively}. 

If executed, the expected outputs are purely numeric arrays residing in the MATLAB workspace:
\begin{itemize}
    \item \textbf{$g$}: A complex double column vector representing the geometric relationship (log-length ratio and angle) of consecutive edges.
    \item \textbf{$\phi$}: An integer column vector acting as a permutation map connecting adjacent edges.
    \item \textbf{$s$}: An integer column vector containing the source vertex indices for every edge.
\end{itemize}

To visualize these results, external MATLAB functions would be required. For example, one could use \texttt{trisurf(T, V(:,1), V(:,2), V(:,3))} to plot the original 3D mesh, or use the \texttt{digraph} and \texttt{plot} functions combined with the source $s$ and target $t$ arrays to visually represent the directed graph of the mesh on a 3D axis.
